\chapter*{はじめる前に}

これは、重さ1.5キログラム、消費電力は薄暗い電球と同じ20ワットでありながら、顔を見分け、言葉を覚え、飛んでくるボールを捕るマシンの本である。つまり、あなたの脳の本だ。エンジニアが見知らぬ回路基板を眺めるように、脳を眺めていこう。どんな部品があり、どうつながっているか。配線をどんな信号が走り、この回路は何を計算しているのか。

数式も生物学の知識もいらない。必要なのは好奇心と、配線や電球やスイッチを思い浮かべる気持ちだけだ。残りは道々組み立てていく。

この本にはもう一つ、完全版がある。方程式、モデル、問題、そして実験への出典つきの版だ。二つの版の章は同じ順序で並び、番号も同じである。ある章が気になったら、その完全版を開いてほしい。どの主張がどこから来たのか、そして自分でどう確かめられるのかが書いてある。本書のサイトでは、ボタン一つで各章を「簡約版」と「完全版」の間で切り替えられる。

ひとつ断っておきたい。脳についてわかっていることは、まだ全体のごく一部にすぎない。どこまでが測定済みで、どこからが推測なのかは、正直に言うつもりだ。これは本書の欠点ではなく、いちばん面白いところである。多くの章は、まだ誰も答えを知らない問いで終わる。
